%==============================================================
% RESULTS SECTION -- journal version.
%
% Every number below is traceable to a file in the project root:
%   ablation rows              ablation_grouped_summary.csv
%   metadata / artefact rows   artefact_control.csv
%   per-split ceilings         split_ceilings.csv
%   text-split performance     textsplit_results.csv, textsplit_proof.csv
%   shift, n=194 controlled    shift_controlled_results.csv
%   shift, n=1308 both ways    big_sealed_results_v2.csv
%   adaptation, 8 seeds        adaptation_final.csv, adaptation_final_seeds.csv
%   article-level screen       article_screen.csv
%
% This section SUPERSEDES the conference version's Section IV-C retraction.
% That retraction rested on an accuracy comparison at one threshold, on 194
% images, with one seed. Subsection E explains what a better measurement says.
%==============================================================

\section{Experimental Results}

\subsection{Per-Modality Ablation}
\label{sec:ablation}

Table~\ref{tab:ablation} reports every configuration under one protocol:
repeated stratified 5-fold \emph{grouped} cross-validation over the
near-duplicate clusters of Section~\ref{sec:dedup}, two repeats, identical
splits and hyperparameters throughout, every transform fit inside the training
fold only, and all figures taken at the conventional $\tau = 0.50$ cut-off.

OCR in isolation reaches $57.35$ percent with a fake-recall of $33.85$. That is
above the 50 percent chance baseline---about three fold-standard-deviations
above it---so the modality is not devoid of signal; but it recovers only a
third of the manipulated class, and adds nothing once CLIP is present. We
remove it, and we do not generalize the result beyond this benchmark:
MultiBanFakeDetect \cite{multibanfake} indicates text-in-image remains a
meaningful vector for Bangla misinformation, and on a corpus where
manipulation is carried by embedded text the ordering could invert.

Every text-only configuration saturates below 78 percent while every
configuration containing CLIP exceeds 95. The gap is not subtle, and
Section~\ref{sec:detectable} explains it: in this split the manipulation is
applied to the pixels and to nothing else.

The comparison that matters for placing this work is against
MultiBanFakeDetect \cite{multibanfake}, the closest prior system in language
and modality scope. We run its feature families---a DenseNet-169 visual
backbone, and that backbone fused with an mBERT text branch---through the
identical pipeline: same folds, same PCA reduction, same classifier, same
cut-off, so that only the backbone differs. DenseNet-169 reaches
$89.28 \pm 1.48$ percent against CLIP's $95.38 \pm 1.33$, a gap of $6.10$
points; with the artefact removed the gap is $6.13$. Adding the mBERT branch
moves DenseNet to $89.65$, which leaves it $5.73$ points behind CLIP alone.

Two things in that comparison are worth separating. The first is the margin
itself, which is large and stable across both encodings, and which we would
expect from a backbone trained contrastively on image--text pairs against one
trained for ImageNet classification. The second is more informative: the mBERT
branch contributes $+0.37$ points to DenseNet, and $+0.45$ with the artefact
removed---the same near-zero contribution the caption pathway makes to CLIP in
Section~\ref{sec:fusion}. Two independently chosen visual backbones give the
same verdict on the text branch, which makes that verdict a property of the
task rather than a quirk of our feature choice.

What this comparison is not is a reimplementation of MultiBanFakeDetect as
published. That system reads the article's own title and body, which our
pipeline never takes as input, so its text branch here encodes the BLIP caption
instead. The comparison therefore isolates the backbone and the fusion
structure under one protocol; it does not establish how the published system
would perform with the inputs it was designed for.

No row is marked best among the four CLIP-containing configurations. Their
total spread is $0.17$ points against a per-fold standard deviation of
$1.0$--$1.4$, and the ordering reverses under an unrelated methodological
change---ungrouped folds put CLIP-only ahead by $0.18$ points, duplicate-aware
folds put the fused configuration ahead by $0.02$. Accuracy does not decide
between them; Section~\ref{sec:fusion} decides on other grounds.

\begin{table}[!tb]
\caption{Baselines, Per-Modality Ablation and the Artefact Control Under One
Protocol: Repeated Stratified 5-Fold \emph{Grouped} Cross-Validation Over
Near-Duplicate Clusters ($2$ Repeats, $10$ Folds), in Percent.}
\label{tab:ablation}
\centering
\setlength{\tabcolsep}{3.5pt}
\footnotesize
\begin{tabular}{lcccc}
\toprule
\textbf{Configuration} & \textbf{Acc.} & \textbf{Prec.$_{f}$} &
\textbf{Rec.$_{f}$} & \textbf{F1$_{f}$} \\
\midrule
\multicolumn{5}{l}{\emph{Reference baselines}} \\
Metadata only (no pixels)$^{\ddagger}$ & 100.00 $\pm$ 0.00 & 100.00 & 100.00 & 100.00 \\
Linear SVM (TF-IDF)     & 73.93 $\pm$ 1.35 & 72.42 & 77.40 & 74.80 \\
XGBoost (TF-IDF + num.) & 77.15 $\pm$ 1.86 & 77.50 & 76.60 & 77.00 \\
DenseNet-169$^{\S}$     & 89.28 $\pm$ 1.48 & 89.10 & 89.50 & 89.30 \\
\quad artefact-controlled & 88.22 $\pm$ 1.57 & 88.49 & 87.90 & 88.19 \\
DenseNet-169 + mBERT$^{\S}$ & 89.65 $\pm$ 1.50 & 89.39 & 90.00 & 89.68 \\
\quad artefact-controlled & 88.67 $\pm$ 1.97 & 88.71 & 88.65 & 88.66 \\
\midrule
\multicolumn{5}{l}{\emph{Modality ablation}} \\
OCR only              & 57.35 $\pm$ 2.34 & 63.97 & 33.85 & 44.22 \\
Caption only          & 74.50 $\pm$ 2.68 & 75.22 & 73.25 & 74.16 \\
Text only             & 76.05 $\pm$ 1.82 & 76.42 & 75.45 & 75.87 \\
Text + CISF           & 77.38 $\pm$ 2.10 & 77.41 & 77.55 & 77.40 \\
\midrule
CLIP only$^{\dagger}$ & 95.38 $\pm$ 1.33 & 94.68 & 96.15 & 95.41 \\
\quad on artefact-controlled images & 94.35 $\pm$ 0.94 & 94.17 & 94.60 & 94.37 \\
Text + CLIP           & 95.28 $\pm$ 1.04 & 95.04 & 95.55 & 95.28 \\
CLIP + CISF           & 95.23 $\pm$ 1.38 & 95.07 & 95.40 & 95.23 \\
Caption + CISF + CLIP & 95.40 $\pm$ 1.09 & 95.13 & 95.70 & 95.41 \\
\bottomrule
\multicolumn{5}{p{0.95\columnwidth}}{\footnotesize Subscript $f$ denotes the
manipulated class; $\pm$ is the spread across folds, not a standard error.
$^{\dagger}$ the selected configuration, on the grounds of
Section~\ref{sec:fusion}. $^{\ddagger}$ width, height, area, aspect ratio, file
size, format and colour mode only, with no image content---see
Section~\ref{sec:artefact}. $^{\S}$ the feature families of
MultiBanFakeDetect \cite{multibanfake} run through this protocol; see the text
for what that comparison does and does not establish.}
\end{tabular}
\end{table}

\subsection{Does the Caption Pathway Contribute?}
\label{sec:fusion}

Accuracy parity does not settle whether the caption pathway is worth keeping: a
feature block can leave accuracy unchanged yet still improve ranking or
calibration, neither of which is visible in Table~\ref{tab:ablation}. We
compare CLIP-only against the fused configuration on four measures in both
regimes, everything else held fixed.

\begin{table}[!tb]
\caption{CLIP-Only Versus the Fused Configuration, Same Hyperparameters and
Folds. AUROC and Brier Are Threshold-Independent; ECE Uses $10$ Bins. Arrows
Mark the Preferred Direction.}
\label{tab:fusion}
\centering
\setlength{\tabcolsep}{4pt}
\footnotesize
\begin{tabular}{llcccc}
\toprule
& \textbf{Model} & \textbf{Acc.}$\uparrow$ & \textbf{AUROC}$\uparrow$ &
\textbf{Brier}$\downarrow$ & \textbf{ECE}$\downarrow$ \\
\midrule
\multirow{2}{*}{\rotatebox{90}{\scriptsize same-d.}}
& CLIP only & 95.38 & 0.9918 & 0.0349 & 3.17 \\
& Fused     & 95.40 & 0.9920 & 0.0342 & 3.05 \\
\midrule
\multirow{2}{*}{\rotatebox{90}{\scriptsize shift}}
& CLIP only & 85.05 & \textbf{0.9759} & \textbf{0.0936} & \textbf{4.06} \\
& Fused     & 85.57 & 0.9732 & 0.1076 & 7.31 \\
\bottomrule
\multicolumn{6}{p{0.94\columnwidth}}{\footnotesize Same-distribution figures
are pooled out-of-fold values over the folds of Table~\ref{tab:ablation};
shift figures are the 194-image partition at $\tau = 0.50$, so both models meet
the same operating point. That partition is easier than the full sealed set:
CLIP-only reaches AUROC $0.976$ on it and $0.941$ on the 1{,}308-image partition
(Fig.~\ref{fig:roc}), so the shift rows support the paired comparison between the
two models, not an absolute level. Both regimes use the original files: see the
caveat at the end of this subsection.}
\end{tabular}
\end{table}

Same-distribution, the two are indistinguishable on every measure. The paired
fold-level differences (fused minus CLIP-only), with the Nadeau--Bengio
correction for fold dependence \cite{nadeau}, are $+0.03$ accuracy points
(95\% CI $[-0.98, +1.03]$, $p = 0.956$), $+0.0002$ AUROC ($p = 0.811$),
$-0.0006$ Brier ($p = 0.709$) and $-0.12$ ECE points ($p = 0.794$); McNemar on
the repeat-averaged out-of-fold predictions, one per image, finds 21 images the
fused model alone classifies correctly against 20 for CLIP-only (exact
$p = 1.000$). Every interval spans zero, and the uncorrected intervals---which
we do not report---were $1.9\times$ narrower.

Under shift the comparison turns against the caption pathway. Its
$+0.5$-point accuracy advantage sits well inside its interval
($[-2.6, +3.6]$), while its calibration is clearly worse: Brier rises from
$0.0936$ to $0.1076$ and expected calibration error from $4.06$ to $7.31$
points---least trustworthy exactly where the threshold has to be chosen. The
pathway costs a BLIP forward pass, an SBERT encoding and a nearest-neighbour
search per image and returns nothing measurable, so we do not claim it as a
contribution and adopt CLIP-only as the recommended configuration.

We repeated the same-distribution comparison under the artefact control of
Section~\ref{sec:artefact}, regenerating the BLIP captions from the re-encoded
images, since the caption pathway reads the image and could carry the artefact
too. With every input taken from the re-encoded files, CLIP-only reaches
$94.35 \pm 0.99$ percent and the fused configuration $94.68 \pm 1.07$: a
difference of $+0.33$ points (Nadeau--Bengio 95 percent CI $[-0.51, +1.16]$),
$-0.0007$ AUROC ($[-0.0023, +0.0008]$), and no difference in Brier score or
expected calibration error; the out-of-fold McNemar test splits $19$ to $17$
($p = 0.87$). The verdict is unchanged. What we did not repeat is the
distribution-shift comparison above, whose calibration gap was measured on the
original files, so that part remains uncontrolled.

\subsection{Dataset Artefacts and What Survives Them}
\label{sec:artefact}

An accuracy above 95 percent on a constructed benchmark deserves suspicion
before it deserves interpretation. In the AdversarialEdit split every one of
the 1{,}000 manipulated images is $1408 \times 768$, while the authentic images
are spread across many resolutions, formats and colour modes; AdversarialGen
carries the same signature in 995 of its 1{,}000 manipulated images. The
classes are separable without reading a single pixel, and the first row of
Table~\ref{tab:ablation} confirms it: a classifier given only width, height,
area, aspect ratio, file size, format and colour mode reaches
$100.00 \pm 0.00$ percent under the same grouped protocol.
Figure~\ref{fig:resolution} shows the fingerprint: the manipulated class sits in
a single column, and the authentic class, spread over 196 distinct sizes in
AdversarialEdit and 157 in AdversarialGen, never uses it.

\begin{figure}[!tb]
\centering
\includegraphics[width=\columnwidth]{fig_j_resolution.png}
\caption{Resolution fingerprint in the two image-attack splits. Manipulated
images are $1408 \times 768$ in $100.0$ percent (a) and $99.5$ percent (b) of
cases; authentic images are spread over many sizes, none of them $1408 \times
768$. Bars show the share of each class at its five most common sizes, plus the
manipulated class's modal size.}
\label{fig:resolution}
\end{figure}

We therefore re-encode every image through one identical pipeline before
feature extraction---resized to $384 \times 384$, converted to RGB, saved as
JPEG at quality 90, EXIF discarded---so that resolution, aspect ratio, format,
colour mode and EXIF are destroyed while content is preserved. Repeating the
CLIP-only configuration on the re-encoded corpus, with the same folds,
hyperparameters and cut-off, gives $94.35 \pm 0.94$ percent against
$95.38 \pm 1.33$ on the originals: a loss of $1.03$ points. The discriminative
signal is in the pixels, not in the storage fingerprint.

The shortcut is not a property of the corpus as a whole, and the benchmark
contains its own control. Table~\ref{tab:shortcut} applies the same
metadata-only classifier to all five adversarial splits. It separates the
classes perfectly in exactly the two splits whose manipulated images are
produced by an image pipeline, and finds nothing in the three that rewrite the
article text and reuse the photograph: AUROC $1.0000$ against $0.3175$ to
$0.4892$. The shortcut therefore tracks the attack type rather than the
benchmark, which is what one would expect if it originates in a generator
emitting at a fixed resolution.

\begin{table}[!tb]
\caption{Metadata-Only Classifier Across All Five Adversarial Splits, Same
Feature Set as Table~\ref{tab:ablation}, 5-Fold Cross-Validation. No Pixel
Content Is Used.}
\label{tab:shortcut}
\centering
\setlength{\tabcolsep}{4pt}
\footnotesize
\begin{tabular}{llcc}
\toprule
\textbf{Split} & \textbf{Manipulation} & \textbf{Acc.} & \textbf{AUROC} \\
\midrule
AdversarialEdit  & image & 100.00 $\pm$ 0.00 & 1.0000 \\
AdversarialGen   & image & 99.35 $\pm$ 1.30 & 1.0000 \\
\midrule
LinguisticAttack & text  & 48.95 $\pm$ 1.61 & 0.4892 \\
RandomSwap       & text  & 39.05 $\pm$ 3.87 & 0.3318 \\
SemanticSwap     & text  & 37.05 $\pm$ 1.54 & 0.3175 \\
\midrule
MultiBanFakeDetect$^{\ddagger}$ & fake-labelled & 77.51 $\pm$ 4.55 & 0.8533 \\
\bottomrule
\multicolumn{4}{p{0.92\columnwidth}}{\footnotesize Values below chance mean the
metadata carries no usable signal and is, if anything, mildly misleading---
consistent with the image reuse across labels documented in
Section~\ref{sec:detectable}. $^{\ddagger}$A second corpus
\cite{multibanfake}: 9{,}597 images, label read from the file name, same
metadata features, plain 5-fold cross-validation.}
\end{tabular}
\end{table}

Two limits on what this establishes. Re-encoding removes the storage artefact,
not a generation-pipeline signature that survives resampling. If every
manipulated image carries the same synthesis fingerprint, a visual model may be
recognising the generator rather than the manipulation---and because BanFakeMM
defines its labels by the construction pipeline, the two are difficult to
separate even in principle. We report the controlled figure throughout and
flag, in each table, which numbers are controlled and which are not.

The control of Table~\ref{tab:shortcut} is internal to one benchmark, so we ran
the same file-metadata classifier on a second Bangla multimodal corpus,
MultiBanFakeDetect \cite{multibanfake}: 9{,}600 images, half labelled fake,
with the label read from the file name. The artefact is present there but far
weaker than in BanFakeMM. Metadata alone reaches $77.5$ percent accuracy
(AUROC $0.853$) under plain five-fold cross-validation, $82.8$ percent
(AUROC $0.902$) with byte-identical files grouped, and $77.1$ percent
(AUROC $0.852$) when trained on the official Train folder and scored on Test,
against $100$ percent in BanFakeMM's image-attack splits and $50$ percent at
chance. One size, $728\times410$, accounts for $39.5$ percent of the
fake-labelled images and for at most $7.2$ percent of any size among the
authentic ones, and $85$ files occur under both labels. We therefore read the
shortcut as a property of constructed Bangla benchmarks that varies in
strength, not as a feature of BanFakeMM alone. Whether the fake-labelled images
in that corpus are manipulated or merely sourced differently we did not
verify, and we use ``fake-labelled'' for that reason.

\subsection{Which Attack Types Are Detectable From Pixels}
\label{sec:detectable}

BanFakeMM provides five adversarial splits. They are not five variations on
one task. We checked each manipulated image against the benchmark's full
authentic corpus (the 7{,}040 real images of the released master file) by
content hash and by perceptual hash; Table~\ref{tab:ceiling} reports the
result.

In AdversarialEdit and AdversarialGen not one manipulated image reproduces an
authentic one. In LinguisticAttack, RandomSwap and SemanticSwap \emph{every}
manipulated image is byte-identical to an authentic image---1{,}000 of 1{,}000,
at perceptual-hash distance zero. Those attacks alter the article text and
leave the photograph untouched. An image-only classifier cannot exceed chance
on them, not because it is weak but because the label is not a function of its
input.

\begin{table}[!tb]
\caption{Per-Split Image-Only Ceiling and Measured Performance. The Ceiling Is
$1$ Minus the Share of Manipulated Images That Reproduce an Authentic Image.
Measured Figures Are the Zero-Shot CLIP-Only Model at $\tau = 0.50$.}
\label{tab:ceiling}
\centering
\setlength{\tabcolsep}{4pt}
\footnotesize
\begin{tabular}{lccccc}
\toprule
\textbf{Split} & \textbf{Manip.} & \textbf{Reproduced} &
\textbf{Ceiling} & \textbf{Acc.} & \textbf{AUROC} \\
\midrule
AdversarialEdit  & 1000 & 0    & 100\% & --- & --- \\
AdversarialGen   & 1000 & 0    & 100\% & 85.1 & --- \\
\midrule
LinguisticAttack & 1000 & 1000 & 0\%   & 49.5 & 0.4832 \\
RandomSwap       & 1000 & 1000 & 0\%   & 50.5 & 0.5165 \\
SemanticSwap     & 1000 & 1000 & 0\%   & 50.5 & 0.5094 \\
\bottomrule
\multicolumn{6}{p{0.95\columnwidth}}{\footnotesize ``Reproduced'' counts
manipulated images byte-identical to an authentic image in the benchmark's
master file; all of them are also at perceptual-hash distance $0$.
AdversarialEdit is the training split, so no held-out figure is quoted.}
\end{tabular}
\end{table}

The measured numbers match the prediction exactly. Accuracy on the three
text-attack splits is $49.5$, $50.5$ and $50.5$ percent, and AUROC is
$0.4832$, $0.5165$ and $0.5094$---indistinguishable from a coin, and
threshold-independent, so no operating point can rescue them. The signature is
visible in the structure as well: fake-recall and real-recall sum to almost
exactly 100 at every threshold we tried, which is what happens when lowering a
cut-off moves images out of one class and into the other at the same rate in
both. Figure~\ref{fig:roc} puts the contrast in one picture: on the
image-attack split the ROC curve climbs steeply, and on the three text-attack
splits it lies on the diagonal.

\begin{figure}[!tb]
\centering
\includegraphics[width=0.95\columnwidth]{fig_j_roc.png}
\caption{ROC curves of the zero-shot CLIP-only model (seed 42, original files)
on the image-attack AdversarialGen split (the 1{,}308-image sealed partition)
and on the three text-attack splits (2{,}000 images each). The text-attack
curves are indistinguishable from chance because every manipulated image in
those splits is byte-identical to an authentic one.}
\label{fig:roc}
\end{figure}

Two further checks confirm the mechanism rather than merely the outcome. Where
a split's manipulated image is the same file as one of its own authentic
images---26, 38 and 131 such pairs respectively---the model assigns the two
\emph{exactly} the same probability, maximum absolute difference
$0.00\mathrm{e}{+}00$. And the mean predicted probability of the manipulated
class ($0.0703$, $0.0678$, $0.0808$) is within noise of the authentic class
($0.0790$, $0.0616$, $0.0739$). The 50 percent is arithmetic, not a deficiency.

The consequence for this paper is a scope statement rather than a result: our
system is image-derived end to end---BLIP caption, OCR and CLIP all read the
image, and the article's own text is never an input---so it is well posed on
two of the five splits and we evaluate it on those two only. The consequence
for the field is sharper: any image-level result reported across all five
BanFakeMM splits is averaging over three conditions in which the task is
undefined.

\subsection{Performance Under Distribution Shift}
\label{sec:shift}

Shift is measured from AdversarialEdit to AdversarialGen, the two splits
Section~\ref{sec:detectable} shows to be well posed. Thresholds are selected on
the disjoint 148-image validation partition under the rule fixed in
Section~\ref{sec:training} and are not reselected afterwards; the sealed
partition is scored once per threshold.

We enlarge the sealed partition well beyond the conference version. The 1{,}350
AdversarialGen images the study had never touched---not training, not
adaptation, not threshold selection---are screened against everything the
models have seen (base training plus the adaptation set) under the full
criteria of Section~\ref{sec:dedup}: perceptual-hash distance at most 2 on
either hash, or CLIP cosine at least $0.98$. That removes 42 images
(21 by pHash, 40 by dHash, 24 by CLIP cosine; 42 in union) and leaves
\textbf{1{,}308 sealed images}, 658 of them manipulated---$6.7\times$ the
194-image partition the conference version used.

\begin{table}[!tb]
\caption{Distribution Shift on the 1{,}308-Image Sealed Partition, Scored Once
Per Threshold, With and Without the Artefact Control. All Values in Percent.}
\label{tab:shift}
\centering
\setlength{\tabcolsep}{3pt}
\footnotesize
\begin{tabular}{llccccc}
\toprule
& \textbf{Setting} & \textbf{$\tau$} & \textbf{Acc.} &
\textbf{Prec.$_{f}$} & \textbf{Rec.$_{f}$} & \textbf{Rec.$_{r}$} \\
\midrule
\multirow{2}{*}{\rotatebox{90}{\scriptsize raw}}
& Zero-shot & 0.09 & 87.2 & 83.1 & 93.5 & 80.8 \\
& Adapted   & 0.12 & 87.8 & 82.0 & 97.1 & 78.5 \\
\midrule
\multirow{2}{*}{\rotatebox{90}{\scriptsize ctrl.}}
& Zero-shot & 0.25 & 84.9 & 86.6 & 82.7 & 87.1 \\
& Adapted   & 0.35 & 88.7 & 89.1 & 88.3 & 89.1 \\
\bottomrule
\multicolumn{7}{p{0.95\columnwidth}}{\footnotesize ``raw'' uses the original
files and carries the artefact of Section~\ref{sec:artefact}; ``ctrl.'' uses
re-encoded images. Bootstrap 95\% CIs on accuracy (5{,}000 resamples), in row
order: $[85.3, 88.9]$, $[86.0, 89.6]$, $[82.9, 86.9]$, $[86.9, 90.4]$.}
\end{tabular}
\end{table}

Two things change under the control, and both matter. The selected thresholds
move from $0.09$ and $0.12$ to $0.25$ and $0.35$: on re-encoded images the
model is far better calibrated and no longer needs an extreme cut-off. And the
apparent benefit of thresholding collapses. On the 194-image partition,
re-encoded, moving from the default cut-off to the validation-selected one
bought $+0.5$ points ($87.6 \rightarrow 88.1$, 95 percent CI
$[-3.6, +4.6]$, McNemar $p = 1.000$), where the uncontrolled measurement had
suggested $+6.1$. \emph{We therefore withdraw the claim, made in the
conference version of this work, that almost the entire shift penalty is a
calibration effect.} It was largely an artefact of the uncontrolled encoding.

\subsection{Does Supervised Adaptation Help?}
\label{sec:adaptation}

The conference version of this work retracted a previously claimed adaptation
gain, on the basis that 300 additional target-style images changed sealed-test
accuracy by $-1.5$ points (95 percent CI $[-4.6, +1.5]$ as reported there; a
bootstrap run to convergence on the same per-image scores gives
$[-5.2, +1.5]$, which is the interval drawn in Fig.~\ref{fig:adaptation}). That
measurement was an accuracy comparison at one fixed threshold, on 194 images,
from a single model fit. It is the most fragile way to measure this effect, and
Table~\ref{tab:adaptation} and Fig.~\ref{fig:adaptation} show it was measuring
the wrong thing.

\begin{table}[!tb]
\caption{Effect of Supervised Adaptation on the 1{,}308-Image Sealed
Partition, Averaged Over Eight Seeds. $\Delta$AUROC Is Threshold-Free; CIs Are
Paired Bootstraps Over Images ($5{,}000$ Resamples).}
\label{tab:adaptation}
\centering
\setlength{\tabcolsep}{4pt}
\footnotesize
\begin{tabular}{lccccc}
\toprule
\textbf{Setting} & \textbf{AUROC$_{0}$} & \textbf{AUROC$_{a}$} &
\textbf{$\Delta$AUROC} & \textbf{95\% CI} & \textbf{$\Delta$Acc.} \\
\midrule
Raw        & 0.9422 & 0.9661 & $+0.0240$ & $[+0.0173, +0.0308]$ & $+1.40 \pm 0.69$ \\
Controlled & 0.9380 & 0.9652 & $+0.0273$ & $[+0.0205, +0.0344]$ & $+4.04 \pm 0.43$ \\
\midrule
\multicolumn{6}{l}{\emph{Article-disjoint subset ($n = 1{,}222$)}} \\
Raw        & 0.9432 & 0.9660 & $+0.0228$ & $[+0.0160, +0.0301]$ & $+1.61 \pm 0.76$ \\
Controlled & 0.9386 & 0.9655 & $+0.0270$ & $[+0.0202, +0.0343]$ & $+4.15 \pm 0.52$ \\
\bottomrule
\multicolumn{6}{p{0.95\columnwidth}}{\footnotesize Subscripts $0$ and $a$
denote the zero-shot and adapted models. $\Delta$Acc. is at the fixed
validation-selected thresholds of Table~\ref{tab:shift}, reported as mean
$\pm$ standard deviation across seeds. Seed-to-seed standard deviation of
$\Delta$AUROC is $0.0018$ (raw) and $0.0017$ (controlled). All four rows use the
same eight seeds, the same training-row order and the same definition of
$\Delta$Acc.}
\end{tabular}
\end{table}

\begin{figure*}[!tb]
\centering
\includegraphics[width=0.97\textwidth]{fig_j_adaptation.png}
\caption{The same intervention measured four ways. (a) Accuracy gain from
adaptation at a fixed threshold, each bar a single fit (seed 42) with a 95
percent bootstrap interval over images: it is negative on the 194-image
partition, indistinguishable from zero on 1{,}308 original files, and positive
on 1{,}308 artefact-controlled images; the eight-seed means are in
Table~\ref{tab:adaptation}. (b) The
threshold-free AUROC gain, with 95 percent intervals, is positive and similar in
both settings. (c) Per-seed AUROC before and after adaptation: it improves in
all 16 runs (eight seeds, two settings).}
\label{fig:adaptation}
\end{figure*}

Measured threshold-free, supervised adaptation improves the ranking by
$+0.0240$ AUROC on the original files and $+0.0273$ on re-encoded ones. Both
confidence intervals exclude zero by a wide margin, and the seed-to-seed
standard deviation---$0.0018$ and $0.0017$---is roughly one fourteenth of the
effect. The result does not depend on the artefact control: it is present, at
almost the same size, with and without it.

Accuracy at a fixed threshold is the unreliable view of the same thing. It
ranges from $+1.40 \pm 0.69$ points to $+4.04 \pm 0.43$ depending only on which
encoding and which operating point one happens to be standing at, and on 194
images with one seed it came out negative. The seed behind the conference
version's single fit (42) gave $+0.7$ on the 1{,}308 original files, the
second-lowest of the eight raw runs, which span $+0.2$ to $+2.6$.

\emph{We therefore reverse the retraction.} Supervised adaptation on 300
target-style images does improve the model, by about $0.025$ AUROC. What the
earlier result established is narrower and still true: the gain is not
reliably visible in accuracy at a fixed operating point, and a study powered
only to detect several points of accuracy will miss it.

\subsection{Article-Level Overlap}
\label{sec:articlescreen}

The screen of Section~\ref{sec:dedup} removes repeated pictures, not repeated
stories. Because every image in this benchmark is named by its source article,
the stronger check is exact: 86 of the 1{,}308 sealed images ($6.6$ percent)
share a source article with a training image. All 86 are authentic;
\emph{none} of the 658 manipulated images does, because AdversarialGen
generates its manipulated images under fresh article identifiers while drawing
its authentic images from the same corpus the training split uses.

Restricting to the 1{,}222 article-disjoint images leaves the conclusions
unchanged, as the lower block of Table~\ref{tab:adaptation} shows:
$\Delta$AUROC moves by $0.0012$ (raw) and $0.0003$ (controlled), both far
inside their confidence intervals, and $\Delta$accuracy by $0.2$ and $0.1$
points, within one seed standard deviation. We report the figure on the full
sealed partition and note that article-level screening does not alter it.
